\documentclass[10pt,a4paper]{amsart}
\usepackage[margin=19mm]{geometry}
\usepackage{amsmath,amssymb,mathtools}
\usepackage[scr=rsfs,cal=euler]{mathalfa}
\usepackage{xfrac}
\usepackage[unicode,hidelinks,bookmarks=false]{hyperref}
\newcommand{\ie}{i.e.\ }
\newcommand{\cf}{cf.\ }
\newcommand{\resp}{respectively\ }
\newcommand{\Subsection}{\subsection}
\providecommand{\nobreakdash}{\nobreak\hbox{-}}
\setlength{\emergencystretch}{2em}
\multlinegap0pt
\usepackage{mathtools}
\usepackage{amssymb}
\usepackage[shortlabels]{enumitem}
\newtheorem{theorem}{Theorem}
\title{Algebraic Geometry over Non-Algebraically Closed Fields:\\ A-Coherent Sheaves over a Ringed Space}
\author{Hamet Seydi, Teylama Miabey, Shurron Farmer}
\date{Source: arXiv:2604.15592v1 (CC BY 4.0)}
\begin{document}
\maketitle

\noindent\emph{Source and licence.} Algebraic Geometry over Non-Algebraically Closed Fields:\\ A-Coherent Sheaves over a Ringed Space. Version arXiv:2604.15592v1 (CC BY 4.0), distributed by arXiv under the Creative Commons Attribution 4.0 International licence, http://creativecommons.org/licenses/by/4.0/, as recorded in the arXiv licence metadata for this exact version and shown on the abstract page https://arxiv.org/abs/2604.15592v1. Full licence notice: \texttt{licenses/arxiv-2604.15592-CC-BY-4.0.txt}.\par\smallskip

\noindent\textit{Editorial scope.} Counted unit: the article's theorem theorem (source0.tex lines 154--162). The article's complete original proof of it is delivered with it (source0.tex lines 164--190, one \texttt{proof} environment of 27 lines). No mathematics is written, repaired or re-derived: the statement and the proof are the article's own lines, in the article's own notation, and no retained line is substituted or re-flowed.  No further source-local context is needed: every cross-reference of every retained line resolves inside the two delivered spans.  Nothing was omitted from any retained span, so the package carries no omission record.  No retained line contains a citation, so the wrapper carries no bibliography and no bibliography entry.  No retained line contains a figure environment, an image-inclusion command or drawing code, so no image is delivered and the package declares no asset dependency.  Editorial formatting only, in addition: an AMS \texttt{amsart} preamble that restates the article's own declarations and macros used by the retained lines, namely the environments \texttt{theorem} on separate counters, together with the standard packages those lines need.  The retained environments print in this document as Theorem 1 in order of appearance, whereas the article prints the same items with the numbering of its own full document.  The complete native source of the article as distributed in the arXiv e-print for this exact version is bundled unchanged as \texttt{sources/198623/source0.tex}.
\begin{theorem}[I.2]
Under the same notations and assumptions as in (I.1), suppose the following conditions hold:
\begin{enumerate}
\item The ring $\Gamma(X, \mathcal{O}_X)$ is Noetherian (resp. $\mathcal{O}_X$ is coherent).
\item $\varphi$ is a flat morphism.
\item The functor $F \mapsto F(X, F)$ from the category of $\mathcal{O}_X$–modules $A$–coherent (resp. $\mathcal{O}_X$–modules coherent) to the category of $\Gamma(X, \mathcal{O}_X)$–modules is exact.
\end{enumerate}
Then the functor $G \mapsto \varphi^* G$ from the category of $\mathcal{O}_Y$–modules coherent (resp. $\mathcal{O}_Y$–modules finitely presented) to the category of $\mathcal{O}_X$–modules $A$–coherent is an equivalence of categories.
\end{theorem}

\begin{proof}
Let $F$ be an $\mathcal{O}_X$–module $A$–coherent. Then there exists an exact sequence
\[
\mathcal{O}_X^m \xrightarrow{u} \mathcal{O}_X^n \longrightarrow F \longrightarrow 0
\]
where $u = \varphi^* v$ for some morphism $v$ of $\mathcal{O}_Y$–modules $\mathcal{O}_Y^m \to \mathcal{O}_Y^n$.  

We deduce that
\[
\ker(u) = \varphi^*(\ker(v)), \quad \mathrm{Im}(u) = \varphi^*(\mathrm{Im}(v)),
\]
since $\varphi$ is flat. As $\ker(u)$ and $\mathrm{Im}(u)$ are coherent if $\mathcal{O}_X$ is coherent, we deduce in both cases that the sequences
\[
0 \to \Gamma(X, \ker(u)) \to \Gamma(X, \mathcal{O}_X^m) \to \Gamma(X, \mathrm{Im}(u)) \to 0
\]
and
\[
0 \to \Gamma(X, \mathrm{Im}(u)) \to \Gamma(X, \mathcal{O}_X^n) \to \Gamma(X, F)
\]
are exact. We thus conclude that the sequence
\[
\Gamma(X, \mathcal{O}_X^m) \to \Gamma(X, \mathcal{O}_X^n) \to \Gamma(X, F)
\]
is exact.

From the proof of (I.1), if $G = F(X, F)^\sim$, then $F \cong \varphi^* G$. Moreover, if $F_1 \to F_2$ is a morphism of $\mathcal{O}_X$–modules $A$–coherent, it is clear that $u = \varphi^* f$ where $\Gamma(X, u)^\sim$, hence the conclusion.
\end{proof}

\end{document}
